With $t=\tan x$ and $t\ne1$, the equation becomes $t+\dfrac{(t+1)}{(1-t)}=1$, giving $t(t-3)=0$. These values satisfy the original domain restrictions. Thus the four points have arguments $0,\theta,\pi,\pi+\theta$, where $\theta=\arctan3$. They form a rectangle with two unit-radius diagonals, and its area is $2\sin\theta$. Since $\sin\theta=\dfrac{3}{\sqrt{10}}$, the area is $\boxed{3\dfrac{\sqrt{10}}{5}}$.
\begin{asy}
texpreamble("\usepackage{amsmath,amssymb}");
size(240);
defaultpen(fontsize(9pt));
pair P=(1,0),Q=(1,3)/sqrt(10);draw(unitcircle);filldraw(P--Q--(-P)--(-Q)--cycle,lightblue,blue);draw(P--(-P),dashed);draw(Q--(-Q),dashed);label("$0$",P,E);label("$\theta$",Q,NE);label("$\pi$",-P,W);label("$\pi+\theta$",-Q,SW);
\end{asy}
